\section{Dodecahedron Labeling}

We fix a combinatorial model of the dodecahedron used throughout the
paper.

\subsection{Vertices and edges}

Let the dodecahedron have $20$ vertices and $30$ edges. We label the
vertices by a set $V = \{0,1,\dots,19\}$.

Edges are unordered pairs $(u,v)$ with $u,v \in V$. The full edge list
is fixed once and for all and determines the combinatorial structure.

\subsection{Faces}

Each face is a $5$-cycle in the edge graph. The $12$ faces are recorded
as cyclically ordered vertex lists
\[
(f_1, f_2, f_3, f_4, f_5).
\]

These face cycles determine the local adjacency relations used in the
transport construction.

\subsection{Flags}

A flag is a triple $(v,e,f)$ consisting of a vertex $v$, an incident
edge $e$, and a face $f$ containing $e$.

The dodecahedron has $120$ flags in total.

\subsection{Transport generators}

The transport system is defined using local moves on flags:
\begin{itemize}
\item vertex move: change the vertex along a fixed edge,
\item edge move: change the edge within a fixed face,
\item face move: change the face across a fixed edge.
\end{itemize}

These generate the flag adjacency structure from which the chamber
graph $\Gsixty$ is constructed.

\subsection{Reproducibility}

All computations in the paper use a fixed implementation of this
labeling, consistent with the edge and face data used in the
computational pipeline described in Appendix~D.
